\documentclass{article}
\usepackage{ctex}
\usepackage{tikz}
\usepackage{amsmath}
\begin{document}

1.2-3

将测量到的电荷量（单位：$ 10^{-19}\rm{C} $）从小到大排序：

$ 6.568,\ 8.204,\ 11.50,\ 13.13,\ 16.48,\ 18.08,\ 19.71,\ 22.89,\ 26.13 $

计算相邻电荷量的差值

$$8.204 - 6.568 = 1.636 \qquad11.50 - 8.204 = 3.296 $$
$$13.13 - 11.50 = 1.63 \qquad16.48 - 13.13 = 3.35 $$
$$18.08 - 16.48 = 1.60 \qquad19.71 - 18.08 = 1.63 $$
$$22.89 - 19.71 = 3.18 \qquad26.13 - 22.89 = 3.24 $$

差值多集中在 $ 1.6 \times 10^{-19}\rm{C} $ 附近，且较大差值（如 $ 3.296 $、$ 3.35 $）约为 $ 1.6 \times 10^{-19}\rm{C} $ 的2倍，符合 $ e $ 的整数倍特征。

将每个电荷量除以对应的整数 $ n $，得到 $ e $ 的候选值：

$$\frac{6.568}{4} = 1.642 \qquad\frac{8.204}{5} = 1.6408 $$
$$\frac{11.50}{7} \approx 1.6429 \qquad\frac{13.13}{8} = 1.6413 $$
$$\frac{16.48}{10} = 1.648 \qquad\frac{18.08}{11} \approx 1.6436 $$
$$\frac{19.71}{12} = 1.6425 \qquad\frac{22.89}{14} \approx 1.635 $$
$$\frac{26.13}{16} \approx 1.6331 $$


取这些候选值的平均值，得到 $ e \approx 1.64 \times 10^{-19}\rm{C} $

最终结论

根据密立根早期实验数据，推得元电荷 $ e \approx 1.64 \times 10^{-19}\rm{C} $


\newpage

1.2-5

%1.2-5 两个点电荷$,q_1=+8.0\mu\rm{C},q_2=-16.0\mu\rm{C}(1\mu\rm{C}=10^{-6}\rm{C})$,相距 $20$ cm.求离它们都是 $20$ cm 处的电场强度$E$.

两个点电荷 $q_1 = +8.0 \mu\rm{C}$ 和 $q_2 = -16.0 \mu\rm{C}$ 相距 $20 \rm{cm}$。

离两者均为 $20\rm{cm}$ 的点 $P$ 与 $q_1$、$q_2$ 构成边长为 $20 \rm{cm}$ 的等边三角形。
$$E = \frac{1}{4\pi\varepsilon_0}\frac{|q|}{r^2}$$
$$\varepsilon_0=8.85\times10^{-12}\rm{C}^2/(\rm{N}\cdot \rm{m}^2)$$
$$q_1 = 8.0 \times 10^{-6} \rm{C} \quad q_2 = -16.0 \times 10^{-6} \rm{C}\quad r = 0.2\rm{m}$$
每个电荷在点 $P$ 产生的电场大小:

$$E_1 = k \frac{|q_1|}{r^2} = 1.8 \times 10^6 \rm{N/C}$$

方向远离 $q_1$(与从 $q_1$ 指向 $P$ 的向量相同)

$$E_2 = k \frac{|q_2|}{r^2} = 3.6 \times 10^6 \rm{N/C}$$

方向指向 $q_2$(与从 $q_2$ 指向 $P$ 的向量相反)

利用几何关系将电场分解为分量并叠加:

设 $q_1$ 在原点 $(-\frac{1}{10},0)$,$q_2$ 在 $(\frac{1}{10}, 0)$,则点 $P$ 坐标为 $(0, \frac{\sqrt{3}}{10})$。

$$E_{1x} = E_1 \cos 60^\circ = 0.9 \times 10^6 \rm{N/C}\quad E_{1y} = E_1 \sin 60^\circ = 9\sqrt{3} \times 10^5 \rm{N/C}$$

$$E_{2x} =  1.8 \times 10^6 \rm{N/C}\quad E_{2y} =-18\sqrt{3} \times 10^5 \rm{N/C}$$

$$E_x = E_{1x} + E_{2x} = 2.7 \times 10^6 \rm{N/C},E_y = E_{1y} + E_{2y} = 9\sqrt{3} \times 10^5 \rm{N/C}$$

$$E = \sqrt{E_x^2 + E_y^2} = 18\sqrt{3}\times 10^5\rm{N/C} \approx 3.12 \times 10^6\rm{N/C}$$

E与 x 轴正方向(从 $q_1$ 指向 $q_2$ 的方向)夹角 $\theta$ 满足:

$$\tan \theta = \frac{E_y}{E_x} = \frac{9\sqrt{3} \times 10^5}{2.7 \times 10^6} = -\frac{\sqrt{3}}{3}$$

即 $\theta = -30^\circ$(顺时针 $30^\circ$),指向右下方。

因此,点 $P$ 处的电场强度大小为 $3.12 \times 10^6 \rm{N/C}$,方向与 $q_1q_2$ 连线成 $30^\circ$ 角指向下方(从 $q_1$ 指向 $q_2$ 为参考,顺时针 $30^\circ$)。


\newpage

1.2-6

%1.2-6 如题图所示, 一电偶极子的电偶极矩 $p = ql,P$ 点到偶极子中心 $O$ 的距离为 $r$, $r$ 与 $l$ 的夹角为 $\theta$. 
%在 $r \gg l$ 时, 求 $P$ 点的电场强度 $E$ 在 $r = OP$ 方向的分量 $E_r$ 和垂直于 $r$ 方向上的分量 $E_\theta$.


设偶极子中心为$O$,场点$P$到$O$的距离为$r$,$r$与$l$的夹角为$\theta$。

$+q$的位置:$O$沿$l$方向外侧$\frac{l}{2}$处

$-q$的位置:$O$沿$l$反方向外侧$\frac{l}{2}$处

点电荷的电场公式:$E=\frac{1}{4\pi\varepsilon_0}\frac{q}{r'^2}e_{r'}$,其中$r'$是场点到点电荷的距离

$e_{r'}$是从点电荷指向场点的单位矢量。

对$+q$,到$P$的距离$r_+ = \left|r-\frac{l}{2}\right|$,由余弦定理:

$$r_+^2 = r^2 + \left(\frac{l}{2}\right)^2 2r\cdot\frac{l}{2}\cos\theta$$

$\frac{l^2}{4r^2}\ll1$,利用二项式展开$(1+x)^{-2}\approx1-2x$($x\ll1$):

$$\frac{1}{r_+^2}\approx\frac{1}{r^2}\left(1+\frac{2l\cos\theta}{r}\right)$$

同理,对$-q$,到$P$的距离$r_= \left|r+\frac{l}{2}\right|$,可得:

$$\frac{1}{r_-^2}\approx\frac{1}{r^2}\left(1-\frac{2l\cos\theta}{r}\right)$$

因为$r\gg l$,$e_{r_+}$和$e_{r_-}$与$e_r$(沿$OP$方向的单位矢量)的夹角极小

设为$\alpha$,几何关系得:

径向分量(沿$e_r$):$\cos\alpha\approx1$(高阶小量忽略)

垂直径向分量(垂直$e_r$,沿$\theta$增大方向):$\sin\alpha\approx\frac{\frac{l}{2}\sin\theta}{r}$

且$e_{r_+}$的垂直分量向上,$e_{r_-}$的垂直分量也向上(因$-q$的电场指向自身,分量方向与$+q$一致)。

1.径向分量$E_r$(沿$e_r$方向)

$+q$的径向电场分量:

$$E_{+r}=\frac{1}{4\pi\varepsilon_0}\frac{q}{r_+^2}\cos\alpha\approx\frac{q}{4\pi\varepsilon_0 r^2}\left(1+\frac{2l\cos\theta}{r}\right)$$

$-q$的径向电场分量(电荷为负,电场指向$-q$):

$$E_{-r}=\frac{1}{4\pi\varepsilon_0}\frac{-q}{r_-^2}\cos\alpha\approx\frac{-q}{4\pi\varepsilon_0 r^2}\left(1-\frac{2l\cos\theta}{r}\right)$$

叠加得总径向分量:
$$E_r = E_{+r}+E_{-r} $$
$$E_r=\frac{q}{4\pi\varepsilon_0 r^2}\left[\left(1+\frac{2l\cos\theta}{r}\right)-\left(1-\frac{2l\cos\theta}{r}\right)\right] =\frac{1}{4\pi\varepsilon_0}\frac{2ql\cos\theta}{r^3}$$

代入$p=ql$,最终径向分量:

$$E_r = e_r \cdot \frac{1}{4\pi\varepsilon_0}\frac{2p\cos\theta}{r^3}$$

2. 垂直径向分量$E_\theta$(垂直$e_r$方向)

$+q$的垂直电场分量:

$$E_{+\theta}=\frac{1}{4\pi\varepsilon_0}\frac{q}{r_+^2}\sin\alpha\approx\frac{q}{4\pi\varepsilon_0 r^2}\cdot\frac{l\sin\theta}{2r}$$

$-q$的垂直电场分量(方向与$+q$一致,叠加增强):

$$E_{-\theta}=\frac{1}{4\pi\varepsilon_0}\frac{-q}{r_-^2}(-\sin\alpha)\approx\frac{q}{4\pi\varepsilon_0 r^2}\cdot\frac{l\sin\theta}{2r}$$

叠加得总垂直分量:

$$E_\theta = E_{+\theta}+E_{-\theta} =\frac{q}{4\pi\varepsilon_0 r^2}\cdot\frac{l\sin\theta}{r}$$

代入$p=ql$,最终垂直分量:

$$E_\theta = e_\theta \cdot \frac{1}{4\pi\varepsilon_0}\frac{p\sin\theta}{r^3}$$





\newpage

1.2-9

%1.2-9 题图所示是另一种电四极子,设 $q$ 和 $l$ 都已知,图中 $P$ 点到电四极子中心 $O$ 的距离为 $x$, $PO$ 与正方形的一对边平行,
%求 $P$ 点的电场强度 $E$. 当 $x \gg l$ 时,求 $P$ 点的电场强度 $E$.

设正方形中心在原点, $ P $ 点在 $ x $ 轴上,距离 $ O $ 为 $ x $,即 $ P = (x, 0) $
$$
A: \left( \frac{l}{2}, \frac{l}{2}, 0 \right),
B: \left( -\frac{l}{2}, \frac{l}{2}, 0 \right),
C: \left( -\frac{l}{2}, -\frac{l}{2}, 0 \right),
D: \left( \frac{l}{2}, -\frac{l}{2}, 0 \right)
$$
$$q_A = +q, \quad q_B = -q, \quad q_C = +q, \quad q_D = -q$$

点 $ P = (x, 0, 0) $,到各电荷的距离:

$$r_A = \sqrt{\left(x - \frac{l}{2}\right)^2 + \frac{l^2}{4}},r_B = \sqrt{\left(x + \frac{l}{2}\right)^2 + \frac{l^2}{4}}$$
$$r_C = \sqrt{\left(x + \frac{l}{2}\right)^2 + \frac{l^2}{4}} = r_B,r_D = \sqrt{\left(x - \frac{l}{2}\right)^2 + \frac{l^2}{4}} = r_A$$

$$\vec{E} = \sum_{i} \frac{k q_i (\vec{r} - \vec{r}_i)}{|\vec{r} - \vec{r}_i|^3}$$

设 $ \vec{r}_P = (x, 0) $。

$$\vec{r}_P - \vec{r}_A = \left( x - \frac{l}{2}, -\frac{l}{2}, 0 \right),|\vec{r}_P - \vec{r}_A| = r_A$$
$$\vec{E}_A = \frac{k q}{r_A^3} \left( x - \frac{l}{2}, -\frac{l}{2}, 0 \right)$$


$$\vec{r}_P - \vec{r}_B = \left( x + \frac{l}{2},-\frac{l}{2},0 \right),r_B = \sqrt{(x+\frac{l}{2})^2 + \frac{l^2}{4}}$$
$$\vec{E}_B = \frac{k (-q)}{r_B^3} \left( x + \frac{l}{2}, -\frac{l}{2}, 0 \right)$$

$$\vec{r}_P - \vec{r}_C = \left( x + \frac{l}{2}, \frac{l}{2}, 0 \right),r_C = r_B,\vec{E}_C = \frac{k q}{r_B^3} \left( x + \frac{l}{2}, \frac{l}{2}, 0 \right)$$

$$\vec{r}_P - \vec{r}_D = \left( x - \frac{l}{2}, \frac{l}{2}, 0 \right),r_D = r_A,\vec{E}_D = \frac{k (-q)}{r_A^3} \left( x - \frac{l}{2}, \frac{l}{2}, 0 \right)$$

$$E_{Ax} = \frac{k q (x - l/2)}{r_A^3}\qquad E_{Bx} = -\frac{k q (x + l/2)}{r_B^3}$$
$$E_{Cx} = \frac{k q (x + l/2)}{r_B^3}\qquad E_{Dx} = -\frac{k q (x - l/2)}{r_A^3}$$

$$ E_{Ax} + E_{Dx} = 0 E_{Bx} + E_{Cx} = 0 $$  

再看 $ E_y $ 分量:

$$E_{Ay} = \frac{k q (-l/2)}{r_A^3}\qquad E_{By} = -\frac{k q (-l/2)}{r_B^3} = +\frac{k q (l/2)}{r_B^3}$$
$$E_{Cy} = \frac{k q (l/2)}{r_B^3}\qquad E_{Dy} = -\frac{k q (l/2)}{r_A^3}$$

所以:
$$E_y = -\frac{k q l/2}{r_A^3} + \frac{k q l/2}{r_B^3} + \frac{k q l/2}{r_B^3} - \frac{k q l/2}{r_A^3}$$
$$E_y = - \frac{k q l}{r_A^3} + \frac{k q l}{r_B^3}$$
$$E_y = k q l \left( \frac{1}{r_B^3} - \frac{1}{r_A^3} \right)$$

$$\vec{E}(P) = \left( 0, \; k q l \left[ \frac{1}{r_B^3} - \frac{1}{r_A^3} \right],\right)$$
其中
$$r_A = \sqrt{(x - l/2)^2 + l^2/4}, \quad r_B = \sqrt{(x + l/2)^2 + l^2/4}.$$

$$k =1/(4\pi\varepsilon_0),p=ql$$

$$E=\frac{p}{4\pi\varepsilon_0}[\frac{1}{\left(x^2-xl+\frac{l^2}{2}\right)^{3/2}}-\frac{1}{\left(x^2+xl+\frac{l^2}{2}\right)^{3/2}}]$$

当 $ x \gg l $ 时,展开 $ r_A^{-3} $ 和 $ r_B^{-3} $。
$$E_y = k q l (r2^{-3} - r1^{-3}) = k q l [ -3l x^{-4} + O(x^{-6}) ] = -3 k q l^2 x^{-4} + \cdots$$
所以,在 $x\gg l$时,电场强度大小为 $E = |E_y| =3 k q l^2 / x^4$,方向沿 y轴负方向(因为系数为负)。

$$E = \frac{3}{4\pi\varepsilon_0} \frac{q l^2}{x^4}$$





\newpage

1.2-10

%1.2-10 均匀带电细棒,设棒长为$2l$,总电荷量为 $q$
%(1)在通过自身端点的垂直面上的场强分布
%(2)在自身的延长线上的场强分布.


 (1)  通过自身端点的垂直面上的场强分布
建立坐标系
设细棒沿 $ x $ 轴放置，端点 $ O $ 在坐标原点，棒的另一端点坐标为 $ (2l, 0) $。
取垂直面($ y-z $ 平面)上任意一点 $ P $，坐标为 $ (0, y) $(由对称性，$ z $ 方向场强分量抵消，只需分析 $ x-y $ 平面)。

取电荷元并分析场强分量

在细棒上取一电荷元 $ dq $，位置坐标为 $ (x, 0) $($ 0\le x\le 2l $)，长度为 $ dx $，则 $ dq = \lambda dx $。

电荷元到 $ P $ 点的距离 $ r = \sqrt{x^2+y^2} $，

由库仑定律，$ dq $ 在 $ P $ 点产生的场强大小为:
$$dE = \frac{1}{4\pi\varepsilon_0} \cdot \frac{dq}{r^2} = \frac{\lambda dx}{4\pi\varepsilon_0(x^2+y^2)}$$

将 $ dE $ 分解为沿 $ x $ 轴负方向的分量 $ dE_x $ 和沿 $ y $ 轴正方向的分量 $ dE_y $，由几何关系 $ \cos\alpha = \frac{x}{r} $、$ \sin\alpha = \frac{y}{r} $，得:
$$dE_x=-dE\cos\alpha=-\frac{\lambda xdx}{4\pi\varepsilon_0(x^2+y^2)^{3/2}}$$
$$dE_y=dE\sin\alpha=\frac{\lambda ydx}{4\pi\varepsilon_0(x^2+y^2)^{3/2}}$$

$$E_x = \int_0^{2l} dE_x = -\frac{\lambda}{4\pi\varepsilon_0} \int_0^{2l} \frac{x dx}{(x^2+y^2)^{3/2}} = -\frac{\lambda}{4\pi\varepsilon_0} \left( \frac{1}{y} - \frac{1}{\sqrt{4l^2+y^2}} \right)$$

$$E_y = \int_0^{2l} dE_y =\frac{\lambda y}{4\pi\varepsilon_0} \int_0^{2l} \frac{dx}{(x^2+y^2)^{3/2}} = \frac{\lambda}{4\pi\varepsilon_0} \cdot \frac{2l}{y\sqrt{4l^2+y^2}}$$

代入 $ \lambda = \frac{q}{2l} $ 化简
$$E_x = -\frac{q}{8\pi\varepsilon_0 l} \left( \frac{1}{y} - \frac{1}{\sqrt{4l^2+y^2}} \right) $$
$$E_y = \frac{q}{4\pi\varepsilon_0 y \sqrt{4l^2+y^2}}$$

 场强的大小与方向
大小:$ E = \sqrt{E_x^2 + E_y^2} $
方向:与 $ y $ 轴夹角 $ \varphi $ 满足 $ \tan\varphi = \left| \frac{E_x}{E_y} \right| $
\newpage
 (2)  自身延长线上的场强分布

同(1)，细棒沿 $ x $ 轴，端点 $ O(0,0) $，另一端 $ (2l,0) $，$ \lambda = \frac{q}{2l} $。

取延长线上任意一点 $ P $，坐标为 $ (x, 0) $

分两种情况讨论:$x>2l $(棒的外侧延长线)、$x<0 $(原点外侧延长线)。

 情况1:$ x>2l $(棒右侧延长线)
取电荷元 $ dq = \lambda dx' $(位置 $ x' \in [0,2l] $)，电荷元到 $ P $ 点的距离 $ r = x - x' $

场强方向沿 $ x $ 正方向，大小 $ dE = \frac{\lambda dx'}{4\pi\varepsilon_0(x - x')^2} $。

积分求总场强:
$$E = \int_0^{2l} dE = \frac{\lambda}{4\pi\varepsilon_0} \int_0^{2l} \frac{dx'}{(x - x')^2} = \frac{\lambda}{4\pi\varepsilon_0} \left. \frac{1}{x - x'} \right|_0^{2l} = \frac{\lambda}{4\pi\varepsilon_0} \left( \frac{1}{x - 2l} - \frac{1}{x} \right)$$

代入 $ \lambda = \frac{q}{2l} $，得:
$$E = \frac{q}{8\pi\varepsilon_0 l} \left( \frac{1}{x - 2l} - \frac{1}{x} \right) \quad (x>2l)$$

 情况2:$ x<0 $(棒左侧延长线)
电荷元到 $ P $ 点的距离 $ r = x' - x $(因 $ x<0 $，$ x' - x > 0 $)，场强方向沿 $ x $ 负方向。

同理积分得:
$$E = -\frac{q}{8\pi\varepsilon_0 l} \left( \frac{1}{-x} - \frac{1}{2l - x} \right) \quad (x<0)$$






\newpage

1.2-12

%1.2-12 一半径为 $R$ 的均匀带电圆环, 总电荷量为 $q$. 
%(1) 求轴线上离环中心 $O$ 为 $x$ 处的场强 $E$
%(2) 画出 $E$-$x$ 曲线
%(3) 轴线上什么地方场强最大,其值是多少


设圆环在 $ xy $ 平面内，圆心在原点，半径为 $ R $，总电荷 $ q $，均匀分布，线电荷密度  
$$\lambda = \frac{q}{2\pi R}.$$
取圆环上一段电荷元 $ dq = \lambda R d\theta $，位于 $ (R\cos\theta, R\sin\theta, 0) $。场点 $ P = (0, 0, x) $
从电荷元到 $ P $ 的矢量:
$$\vec{r} = (0 - R\cos\theta, 0 - R\sin\theta, z - 0) = (-R\cos\theta, -R\sin\theta, z).$$
$$r = \sqrt{R^2 + z^2}.$$
由对称性，圆环在 $ P $ 点的电场只有 $ z $ 方向分量，$ x, y $ 分量总和为 0。
电荷元产生的电场在 $ z $ 方向的分量:
$$dE_z = \frac{1}{4\pi\varepsilon_0} \frac{dq}{r^2} \cdot \frac{z}{r}= \frac{1}{4\pi\varepsilon_0} \frac{\lambda R d\theta}{(R^2 + z^2)^{3/2}} \, z.$$
$$E_z = \frac{1}{4\pi\varepsilon_0} \frac{\lambda R z}{(R^2 + z^2)^{3/2}} \int_0^{2\pi} d\theta= \frac{1}{4\pi\varepsilon_0} \frac{\lambda R z}{(R^2 + z^2)^{3/2}} \cdot 2\pi.$$
代入 $ \lambda = \frac{q}{2\pi R} $:
$$E_z = \frac{1}{4\pi\varepsilon_0} \frac{q z}{(R^2 + z^2)^{3/2}}.$$
所以，轴线上距离环心 $ x $ 处的场强(沿轴线方向):
$$E(x) = \frac{1}{4\pi\varepsilon_0} \frac{q x}{(R^2 + x^2)^{3/2}}$$
方向沿轴线，$ x>0 $ 时向外，$ x<0 $ 时向环心。

 (2) 画出 $ E $-$ x $ 曲线

函数 $ E(x) = k \frac{q x}{(R^2 + x^2)^{3/2}} $，其中 $ k = \frac{1}{4\pi\varepsilon_0} $。


\begin{tikzpicture}[scale=1.5, domain=-3:3, samples=200]

    % 坐标轴
    \draw[->] (-3.2,0) -- (3.2,0) node[right] {$x$};
    \draw[->] (0,-1.5) -- (0,1.5) node[above] {$E(x)$};

    % 函数公式标注
    \node[above right] at (2,1) {$ E(x) = k \frac{q x}{(R^2 + x^2)^{3/2}} $};

    % 画出函数曲线
    \draw[thick, blue] plot (\x, {3*\x / ( (1+(\x)^2)^(3/2) )});

    % 真正的极值点 x = ±1/sqrt(2) ≈ ±0.707
    \pgfmathsetmacro{\ex}{1/sqrt(2)}
    \pgfmathsetmacro{\ey}{3*\ex / ((1+(\ex)^2)^(3/2))} % ey = 2/sqrt(3) ≈ 1.1547
    \draw[dashed] (\ex,0) -- (\ex,\ey);
    \draw[dashed] (-\ex,0) -- (-\ex,-\ey);

    % 原点标注
    \node[below right] at (0,0) {$O$};

\end{tikzpicture}



(3) 考虑 $ x \ge 0 $ 部分，令
$$f(x) = \frac{x}{(R^2 + x^2)^{3/2}}.$$
$$f'(x) = \frac{(R^2 + x^2)^{3/2} - x \cdot \frac{3}{2}(R^2 + x^2)^{1/2} \cdot 2x}{(R^2 + x^2)^3}.$$
化简分子(提取 $ (R^2 + x^2)^{1/2} $):
$$(R^2 + x^2)^{1/2} \left[ (R^2 + x^2) - 3x^2 \right] = (R^2 + x^2)^{1/2} (R^2 - 2x^2).$$
所以:
$$f'(x) = \frac{(R^2 + x^2)^{1/2} (R^2 - 2x^2)}{(R^2 + x^2)^3}= \frac{R^2 - 2x^2}{(R^2 + x^2)^{5/2}}.$$
令 $ f'(x) = 0 $ 得 $ R^2 - 2x^2 = 0 $，即
$$x = \pm \frac{R}{\sqrt{2}}.$$
当 $ 0 < x < R/\sqrt{2} $ 时 $ f'(x) > 0 $，当 $ x > R/\sqrt{2} $ 时 $ f'(x) < 0 $

所以 $ x = R/\sqrt{2} $ 处 $ f $ 取极大值。

$$E_{\max} = \frac{1}{4\pi\varepsilon_0} \frac{q \cdot \frac{R}{\sqrt{2}}}{\left(R^2 + \frac{R^2}{2}\right)^{3/2}}= \frac{1}{4\pi\varepsilon_0} \frac{q R/\sqrt{2}}{\left(\frac{3R^2}{2}\right)^{3/2}}.$$
$$\left(\frac{3R^2}{2}\right)^{3/2} = \left(\frac{3}{2}\right)^{3/2} R^3.$$
$$E_{\max} = \frac{1}{4\pi\varepsilon_0} \frac{q R}{\sqrt{2}} \cdot \frac{1}{\left(\frac{3}{2}\right)^{3/2} R^3}= \frac{1}{4\pi\varepsilon_0} \frac{q}{R^2} \cdot \frac{1}{\sqrt{2} \cdot \left(\frac{3}{2}\right)^{3/2}}.$$
$$\sqrt{2} \cdot \left(\frac{3}{2}\right)^{3/2} = \sqrt{2} \cdot \frac{3\sqrt{3}}{2\sqrt{2}} \cdot \frac{1}{\sqrt{2}}$$
$$E_{\max} =\frac{q}{6\sqrt{3}\pi\varepsilon_0R^2}.$$




\newpage

1.2-13

%1.2-13 半径为$R$的圆面上均匀带电,电荷的面密度为$\sigma_e.$
%(1) 求 轴 线上离圆心的坐标为 $x$处的场强；
%(2)在保持$\sigma_e$不变的情况下,当$R\to0$和$R\to\infty$时结果各如何？
%(3)在保持总电荷$Q=\pi R^2\sigma_e$不变的情况下,当$R\to0$和$R\to\infty$时结果各如何？

设圆面半径为$R$，电荷面密度$\sigma_e$，轴线为$x$轴，圆心在坐标原点$O$，场点$P$坐标为$(x, 0)$。

在圆面上取一半径为$r$、宽度为$dr$的同心细圆环，其面积 $dS = 2\pi r dr$，带电量

$$dq = \sigma_e dS = 2\pi\sigma_e r dr$$

已知半径为$r$、带电量为$dq$的细圆环，在轴上距离圆心$x$处的场强大小为

$$dE = \frac{1}{4\pi\varepsilon_0} \cdot \frac{x dq}{(x^2 + r^2)^{3/2}}$$

场强方向沿$x$轴($\sigma_e>0$时背离圆心，$\sigma_e<0$时指向圆心)。

将$dq$代入，积分区间$r \in [0, R]$:
$$E= \int_0^R dE = \int_0^R \frac{1}{4\pi\varepsilon_0} \cdot \frac{x \cdot 2\pi\sigma_e r dr}{(x^2 + r^2)^{3/2}} \frac{\sigma_e x}{2\varepsilon_0} \int_0^R \frac{r dr}{(x^2 + r^2)^{3/2}}$$

令$u = x^2 + r^2$，则$du = 2r dr$，积分上下限变为$u=x^2 \to u=x^2+R^2$，积分式化简为

$$\int \frac{r dr}{(x^2 + r^2)^{3/2}} = \frac{1}{2}\int u^{-3/2} du = -u^{-1/2} + C$$

代入计算得
$$E = \frac{\sigma_e x}{2\varepsilon_0} \left[ \frac{1}{|x|} - \frac{1}{\sqrt{x^2 + R^2}} \right]$$

考虑场强方向，最终轴上场强表达式为

$$E = \frac{\sigma_e}{2\varepsilon_0} \left[ 1 - \frac{|x|}{\sqrt{x^2 + R^2}} \right] e_x$$


 (2) 保持$\sigma_e$不变，分析$R\to0$和$R\to\infty$的极限
当$R\to0$时

$$\frac{|x|}{\sqrt{x^2 + R^2}} \to \frac{|x|}{|x|} = 1$$

代入场强公式得

$$E \to \frac{\sigma_e}{2\varepsilon_0} \left( 1 - 1 \right) = 0$$

当$R\to\infty$时

$$\frac{|x|}{\sqrt{x^2 + R^2}} = \frac{|x|}{R\sqrt{1 + \frac{x^2}{R^2}}} \to 0$$

代入场强公式得

$$E \to \frac{\sigma_e}{2\varepsilon_0}$$

 (3) 保持总电荷$Q=\pi R^2\sigma_e$不变，分析$R\to0$和$R\to\infty$的极限

由$Q=\pi R^2\sigma_e$得$\sigma_e = \frac{Q}{\pi R^2}$，代入轴上场强公式:

$$E = \frac{Q}{2\pi\varepsilon_0 R^2} \left[ 1 - \frac{|x|}{\sqrt{x^2 + R^2}} \right]$$

当$R\to0$时

对括号内部分做泰勒展开($R\ll|x|$):

$$\frac{|x|}{\sqrt{x^2 + R^2}} = \frac{1}{\sqrt{1 + \frac{R^2}{x^2}}} \approx 1 - \frac{R^2}{2x^2}$$

$$E \approx \frac{Q}{2\pi\varepsilon_0 R^2} \cdot \frac{R^2}{2x^2} = \frac{Q}{4\pi\varepsilon_0 x^2}$$

当$R\to\infty$时

$$\frac{|x|}{\sqrt{x^2 + R^2}} \to 0$$

但$\sigma_e = \frac{Q}{\pi R^2}\to0$，代入场强公式得

$$E = \frac{Q}{2\pi\varepsilon_0 R^2} \to 0$$




\end{document}